\section{Considerações finais}
Esta entrega mantém a direção estratégica construída nas versões anteriores, mas fortalece o rigor da tomada de decisão. A PetCuida deve priorizar validação territorial, densidade bilateral e economia unitária antes de escalar. O setor pet é amplo, mas a oportunidade específica precisa ser demonstrada por comportamento de uso, oferta local e pagamento efetivo. A expansão deve ocorrer por evidências e gatilhos, não por metas isoladas de downloads ou usuários cadastrados.

A decisão central é executar um piloto assistido de oito semanas, medir ativação, retenção, CAC, conversão e margem real, e só então atualizar o modelo financeiro. A escalabilidade dependerá de automação de processos estáveis, infraestrutura segura, receita recorrente contratada e redução ou financiamento explícito do custo dos créditos solidários.